
\medskip
\phantomsection\label{prop:sga3-VIB-3.11}%
\noindent\textbf{Proposition 3.11.---}%
\footnote{\label{note:sga3-VIB-26}Editor's note: compare
VI\(_A\), \SGARefLink{par:sga3-VIB-5.6}{5.6}.}
\emph{Let \(u:G\to H\) be a morphism between \(S\)-groups locally
of finite presentation. If \(u\) is flat, the induced map}
\[
  u^0:G^0\longrightarrow H^0
\]
\emph{is surjective. The converse holds if \(G\) is flat over \(S\)
and \(H\) has reduced fibers.}

Suppose that \(u\) is flat. Then, for every \(s\in S\), \(u_s\) is
flat and locally of finite presentation, hence open
(EGA~IV\(_2\), 2.4.6). Therefore
\[
  u_s^0:\underline{G_s^0}\longrightarrow\underline{H_s^0}
\]
is surjective, by \SGARefLink{cor:sga3-VIB-1.3.2}{1.3.2}.%
\footnote{\label{note:sga3-VIB-27}Editor's note: we have shortened
the original by referring here to \SGARefLink{cor:sga3-VIB-1.3.2}{1.3.2}. We have also simplified
the argument below by omitting a reference to the generic flatness
theorem (EGA~IV\(_2\), 6.9.1), which is not needed here.}
Thus the map
\[
  u^0:G^0\longrightarrow H^0
\]
is surjective.

Conversely, suppose that \(u^0:G^0\to H^0\) is surjective, that
\(G\) is flat over \(S\), and that \(H\) has reduced fibers. Then,
for every \(s\in S\), the morphism
\[
  u_s^0:\underline{G_s^0}\longrightarrow\underline{H_s^0}
\]
is surjective, and hence is flat at every point lying above the
generic point \(\eta\) of \(\underline{H_s^0}\), since
\[
  \mathcal O_{\underline{H_s^0},\eta}
\]
is a field. It follows from \hyperref[prop:sga3-VIB-1.3]{1.3} that \(u_s\) is flat. Therefore
\(u\) is flat, by the fiberwise flatness criterion
(EGA~IV\(_3\), 11.3.11).

\section*{4. Dimension of the fibers of groups locally of finite presentation}
\phantomsection\label{sec:sga3-VIB-4}%
\addcontentsline{toc}{section}{4. Dimension of the fibers of groups locally of finite presentation}

\medskip
\phantomsection\label{prop:sga3-VIB-4.1}%
\noindent\textbf{Proposition 4.1.---}
\emph{Let \(G\) be an \(S\)-scheme locally of finite type, equipped
with an \(S\)-section \(\epsilon\), and suppose that, for every
\(s\in S\),}
\[
  \dim G_s=\dim_{\epsilon(s)}G_s
\]
\emph{\textup{(as is the case when \(G\) is an \(S\)-group;
\hyperref[lem:sga3-VIB-1.5]{1.5})}.}
\begin{enumerate}
\item[\textup{\SGAInlineRefTarget{prop:sga3-VIB-4.1:item:i}(i)}] \emph{The function}
  \[
    s\longmapsto\dim G_s
  \]
  \emph{is upper semicontinuous on \(S\).}
\item[\textup{\SGAInlineRefTarget{prop:sga3-VIB-4.1:item:ii}(ii)}] \emph{If, moreover, \(G\) is locally of finite
  presentation over \(S\), this function is locally constructible.}
\end{enumerate}

Let
\[
  \pi:G\longrightarrow S
\]
be the structural morphism. Chevalley's semicontinuity theorem
(EGA~IV\(_3\), 13.1.3) asserts that the function
\[
  x\longmapsto\dim_x\pi^{-1}\bigl(\pi(x)\bigr)
\]
is upper semicontinuous on \(G\). For every \(s\in S\), however,
\[
  \dim G_s
    =\dim\pi^{-1}(s)
    =\dim_{\epsilon(s)}
       \pi^{-1}\bigl(\pi(\epsilon(s))\bigr).
\]
Since the map \(s\mapsto\epsilon(s)\) is continuous on \(S\), the
composite function
\[
  s\longmapsto\dim G_s
\]
is upper semicontinuous on \(S\).

Now suppose that \(G\) is locally of finite presentation over
\(S\). Arguing as above, to show that \(s\mapsto\dim G_s\) is
locally constructible, it is enough to show that the function
\[
  x\longmapsto\dim_x\pi^{-1}\bigl(\pi(x)\bigr)
\]
is locally constructible on \(G\), which follows from
EGA~IV\(_3\), 9.9.1.
